\scanpage{142}
Pour expliciter la condition (iii), il reste donc à expliciter les tores
quotients \inserted{\uline{compacts}} de \deleted{$\unsure{T_L}$} $V_L$, et
prouver que ce sont les quotients de $V_L/V_{L_0}$.

\deleted{$V_{L_0}$ n'a manifestement pas de \inserted{tore} quotient compact
\inserted{sauf le gr. unité}, donc le plus grand quotient compact de $V_L$
est majoré par $V_L/V_{L_0}$.}

\deleted{\unsure{Pour voir} Il reste à prouver que ce dernier est compact.
C'est trivial si $L=L_0$. Sinon, (faisant $L=K$ pour simplifier les notations),
considérons \unsure{le premier fait} i.e. à prouver
\unsure{que} $V_L/V_{L_0}$
\[\unsure{L_{\mathbb R}}\simeq\mathbb R^{\unsure{r_1}}\qquad L_{0\mathbb R}=\]}

\uline{Corollaire 1} Supposons $K$ weilien, soit $K_0=K\cap\mathbb R$,
\deleted{et} Alors $V_K/V_{K_0}$ est compact, et $V_{K_0}$ n'a d'autre
\inserted{tore} quotient compact que le gr unité.
\deleted{L'image de $E_K$ dans}
De plus $E_K\cap V_{K_0}(\mathbb Q)$ est d'indice fini dans $E_K$;
\inserted{i.e. l'image de $E_{1K}$ dans $V_K/V_{K_0}$ est finie.}

Il suffit \inserted{(pour la 1ère assertion)} de prouver l'énoncé analogue
après \unsure{extension} de la base $\mathbb Q-\mathbb R$.
Alors $K_{0\mathbb R}\simeq\mathbb R^r$, \deleted{\unsure{et est un}}
\inserted{\unsure{dans $V_{L_0}$ et à fortiori}}
$\unsure{V_{L_0}}$ est un \inserted{tore} splitté, donc tout quotient est splitté;
donc compact seulement s'il est \deleted{\unsure{le}} nul.
D'autre part $V_K/V_{K_0}$ est compact: c'est trivial si $K=K_0$,
sinon on aura $K_{\mathbb R}\simeq\mathbb C^{\unsure{*}r}$, et
\[
V_K/V_{K_0}\simeq T_K/T_{K_0}\simeq
\deleted{\prod_{\unsure{r}}}
\bigl[(\prod_{\mathbb C/\mathbb R}G_m)/G_m\bigr]^r\simeq U^r
\]
est compact.

La dernière assertion signifie \deleted{\unsure{signifie}} que $E_{K_0}$ est d'indice
fini dans $E_K$. Or les rangs sont respectivmt, si $[K_0:K]=r$,
$r-1$ et $r-1$, O.K.

$(\mathrm{iv})\Longrightarrow(\mathrm{i})$ Résulte immédiatement du fait que
$\unsure{N_{L/K}|N_{K/L}}$ applique $E_{1K}$ dans $E_{1L}$, et de la dernière
assertion du corollaire.
\margin{141}
